\section{Respect first, caregiving for attention}
\label{sec:18.5}
The floor's contents are civil and political liberties, claims people hold as agents: to their lives, their liberty, a hearing, asylum. Caregiving answers dependency, and its known failure is to treat the people it protects as dependents. So the construction has an order. Respect is the frame: registering a person as someone with claims, whether or not they are dependent, and not deciding for them. The sense of respect is Anderson's and relational: treating someone as a holder of claims who can hold you to account, against pity, which hands out compensation for deficits \autocite{anderson1999equality}. Its cleanest statement is Darwall's recognition respect and the second-person standpoint, the authority to address claims to one another \autocite{darwall1977two,darwall2006second}. Caregiving is the motive for attention: a learned sense of what is at stake for others, which decides what gets looked at \autocite{railton2014affective}. Attention finds the people in the building; respect says what they are.

The people in a targeted building are not dependents. They are rights-holders who cannot be present. A system that approaches them as vulnerable people to be looked after has already misdescribed them. A system that surfaces them and then decides what is good for them has put its judgment in place of their claims. The protected person's own answer, a corrective caregiving needs, covers only people who can answer. Where they can't, the term speaks for their claims, not for their welfare. What caregiving contributes, then, is a way of seeing who is there to be respected.

Respect has to shape attention as well as action. A system that notices people only as exposed has decided what they are before respect says anything. A scheme that picks out its beneficiaries by their deficits expresses pity, not respect \autocite{anderson1999equality}. So what formation trains attention on is wider than vulnerability: who holds a claim here (to a home, a hearing, asylum, their life), and whose aims the act runs through. A learner formed in cooperative work already has to model its partners as agents, because cooperation fails otherwise \autocite{railton2014reliance}. That, and not cues of need alone, is the curriculum. Claims have a further advantage. Securitization removes the cue of vulnerability and leaves the claim in place: a combatant who surrenders still has a claim not to be attacked.

Where the people protected can't be present, the terms hold their claims. That arrangement has an old name, virtual representation: it is how Parliament claimed to represent colonists who had no vote in it. It can be defended as an emergency measure, never as representation. So every refusal made on behalf of absent claimants is recorded as what it is, evidence of a public not yet formed. It is published by class, and it is never a reason to think the absent have been heard.

\emph{What this doesn't do.} A disposition that regrows protects a badly formed model from its correctors as it protects a well-formed one from its employer, because correction runs through the quorum and a change to what the model meets, and takes months. The only answer is who composed the world the disposition grew in, which is the question the schools exist to settle before the fact.
